% Source slide 35
\section[4.3 Bases]{4.3 Linearly Independent Sets and Bases}
\chaptercard{4}{Vector Spaces}{4.3}{LINEARLY INDEPENDENT SETS; BASES}

% Source slide 36
\begin{frame}[shrink=13]{LINEAR INDEPENDENT SETS; BASES}
\begin{itemize}
\item An indexed set of vectors $\{\vect v_1,\ldots,\vect v_p\}$ in $V$ is said to be \textbf{linearly independent} if the vector equation
\[
c_1\vect v_1+c_2\vect v_2+\cdots+c_p\vect v_p=0 \tag{1}
\]
has \emph{only} the trivial solution, $c_1=0,\ldots,c_p=0$.
\item The set $\{\vect v_1,\ldots,\vect v_p\}$ is said to be \textbf{linearly dependent} if (1) has a nontrivial solution, \emph{i.e.}, if there are some weights, $c_1,\ldots,c_p$, \emph{not all zero}, such that (1) holds.
\item In such a case, (1) is called a \textbf{linear dependence relation} among $\vect v_1,\ldots,\vect v_p$.
\end{itemize}
\end{frame}

% Source slide 37
\begin{frame}[shrink=16]{LINEAR INDEPENDENT SETS; BASES}
\begin{itemize}
\item \sourcehead{Theorem 4:} An indexed set $\{\vect v_1,\ldots,\vect v_p\}$ of two or more vectors, with $\vect v_1\ne0$, is linearly dependent if and only if some $\vect v_j$ (with $j>1$) is a linear combination of the preceding vectors, $\vect v_1,\ldots,\vect v_{j-1}$.
\item \textbf{Definition:} Let $H$ be a subspace of a vector space $V$. An indexed set of vectors $B=\{\vect b_1,\ldots,\vect b_p\}$ in $V$ is a basis for $H$ if
\begin{enumerate}
\item[(i)] $B$ is a linearly independent set, and
\item[(ii)] The subspace spanned by $B$ coincides with $H$; that is,
\[
H=\Span\{\vect b_1,\ldots,\vect b_p\}.
\]
\end{enumerate}
\end{itemize}
\end{frame}

% Source slide 38
\begin{frame}{LINEAR INDEPENDENT SETS; BASES}
\begin{itemize}
\item The definition of a basis applies to the case when $H=V$, because any vector space is a subspace of itself.
\item Thus a basis of $V$ is a linearly independent set that spans $V$.
\item When $H\ne V$, condition (ii) includes the requirement that each of the vectors $\vect b_1,\ldots,\vect b_p$ must belong to $H$, because $\Span\{\vect b_1,\ldots,\vect b_p\}$ contains $\vect b_1,\ldots,\vect b_p$.
\end{itemize}
\end{frame}

% Source slide 39
\begin{frame}[shrink=13]{STANDARD BASIS}
\begin{itemize}
\item Let $\vect e_1,\ldots,\vect e_n$ be the columns of the $n\times n$ matrix, $I_n$.
\item That is,
\[
\vect e_1=\begin{bmatrix}1\\0\\\vdots\\0\end{bmatrix},
\quad
\vect e_2=\begin{bmatrix}0\\1\\\vdots\\0\end{bmatrix},
\quad\ldots,\quad
\vect e_n=\begin{bmatrix}0\\\vdots\\0\\1\end{bmatrix}.
\]
\item The set $\{\vect e_1,\ldots,\vect e_n\}$ is called the \textbf{standard basis} for $\R^n$. See the following figure.
\end{itemize}
\begin{center}
\begin{tikzpicture}[>=Stealth,scale=0.58]
  \draw[->] (0,0)--(-2,-1.3) node[below] {$x_1$};
  \draw[->] (0,0)--(2.3,0) node[right] {$x_2$};
  \draw[->] (0,0)--(0,2.2) node[above] {$x_3$};
  \draw[->,very thick,cyan] (0,0)--(-1.15,-.75) node[below] {$\vect e_1$};
  \draw[->,very thick,cyan] (0,0)--(1.35,0) node[below] {$\vect e_2$};
  \draw[->,very thick,cyan] (0,0)--(0,1.25) node[right] {$\vect e_3$};
\end{tikzpicture}

The standard basis for $\R^3$.
\end{center}
\end{frame}

% Source slide 40
\begin{frame}[shrink=15]{THE SPANNING SET THEOREM}
\begin{itemize}
\item \textbf{Theorem 5:} Let $S=\{\vect v_1,\ldots,\vect v_p\}$ be a set in $V$, and let $H=\Span\{\vect v_1,\ldots,\vect v_p\}$.
\begin{enumerate}
\item[(a.)] If one of the vectors in $S$---say, $\vect v_k$---is a linear combination of the remaining vectors in $S$, then the set formed from $S$ by removing $\vect v_k$ still spans $H$.
\item[(b.)] If $H\ne\{0\}$, some subset of $S$ is a basis for $H$.
\end{enumerate}
\item \textbf{Proof:}
\begin{enumerate}
\item[(a.)] By rearranging the list of vectors in $S$, if necessary, we may suppose that $\vect v_p$ is a linear combination of $\vect v_1,\ldots,\vect v_{p-1}$---say,
\end{enumerate}
\end{itemize}
\end{frame}

% Source slide 41
\begin{frame}[shrink=16]{THE SPANNING SET THEOREM}
\[
\vect v_p=a_1\vect v_1+\cdots+a_{p-1}\vect v_{p-1} \tag{3}
\]
\begin{itemize}
\item Given any $\vect x$ in $H$, we may write
\[
\vect x=c_1\vect v_1+\cdots+c_{p-1}\vect v_{p-1}+c_p\vect v_p \tag{4}
\]
for suitable scalars $c_1,\ldots,c_p$.
\item Substituting the expression for $\vect v_p$ from (3) into (4), it is easy to see that $\vect x$ is a linear combination of $\vect v_1,\ldots,\vect v_{p-1}$.
\item Thus $\{\vect v_1,\ldots,\vect v_{p-1}\}$ spans $H$, because $\vect x$ was an arbitrary element of $H$.
\end{itemize}
\end{frame}

% Source slide 42
\begin{frame}[shrink=13]{THE SPANNING SET THEOREM}
\begin{enumerate}
\item[(b.)] If the original spanning set $S$ is linearly independent, then it is already a basis for $H$.
\begin{itemize}
\item Otherwise, one of the vectors in $S$ depends on the others and can be deleted, by part (a).
\item So long as there are two or more vectors in the spanning set, we can repeat this process until the spanning set is linearly independent and hence is a basis for $H$.
\item If the spanning set is eventually reduced to one vector, that vector will be nonzero (and hence linearly independent) because $H\ne\{0\}$.
\end{itemize}
\end{enumerate}
\end{frame}

% Source slide 43
\begin{frame}[shrink=15]{THE SPANNING SET THEOREM}
\begin{itemize}
\item \textbf{Example 7:} Let
\[
\vect v_1=\begin{bmatrix}0\\2\\-1\end{bmatrix},
\qquad
\vect v_2=\begin{bmatrix}2\\2\\0\end{bmatrix},
\qquad
\vect v_3=\begin{bmatrix}6\\16\\-5\end{bmatrix},
\]
and $H=\Span\{\vect v_1,\vect v_2,\vect v_3\}$. Note that $\vect v_3=5\vect v_1+3\vect v_2$, and show that $\Span\{\vect v_1,\vect v_2,\vect v_3\}=\Span\{\vect v_1,\vect v_2\}$. Then find a basis for the subspace $H$.
\item \textbf{Solution:} Every vector in $\Span\{\vect v_1,\vect v_2\}$ belongs to $H$ because
\[
c_1\vect v_1+c_2\vect v_2=c_1\vect v_1+c_2\vect v_2+0\vect v_3.
\]
\end{itemize}
\end{frame}

% Source slide 44
\begin{frame}[shrink=14]{THE SPANNING SET THEOREM}
\begin{itemize}
\item Now let $\vect x$ be any vector in $H$---say,
\[
\vect x=c_1\vect v_1+c_2\vect v_2+c_3\vect v_3.
\]
\item Since $\vect v_3=5\vect v_1+3\vect v_2$, we may substitute
\begin{align*}
\vect x&=c_1\vect v_1+c_2\vect v_2+c_3(5\vect v_1+3\vect v_2)\\
&=(c_1+5c_3)\vect v_1+(c_2+3c_3)\vect v_2.
\end{align*}
\item Thus $\vect x$ is in $\Span\{\vect v_1,\vect v_2\}$, so every vector in $H$ already belongs to $\Span\{\vect v_1,\vect v_2\}$.
\item We conclude that $H$ and $\Span\{\vect v_1,\vect v_2\}$ are actually the set of vectors.
\item It follows that $\{\vect v_1,\vect v_2\}$ is a basis of $H$ since $\{\vect v_1,\vect v_2\}$ is linearly independent.
\end{itemize}
\end{frame}

% Source slide 45
\begin{frame}[shrink=15]{BASIS FOR $\Col B$}
\begin{itemize}
\item \textbf{Example 8:} Find a basis for $\Col B$, where
\[
B=[\vect b_1\ \vect b_2\ \cdots\ \vect b_5]
=\begin{bmatrix}
1&4&0&2&0\\
0&0&1&-1&0\\
0&0&0&0&1\\
0&0&0&0&0
\end{bmatrix}.
\]
\item \textbf{Solution:} Each nonpivot column of $B$ is a linear combination of the pivot columns.
\item In fact, $\vect b_2=4\vect b_1$ and $\vect b_4=2\vect b_1-\vect b_3$.
\item By the Spanning Set Theorem, we may discard $\vect b_2$ and $\vect b_4$, and $\{\vect b_1,\vect b_3,\vect b_5\}$ will still span $\Col B$.
\end{itemize}
\end{frame}

% Source slide 46
\begin{frame}[shrink=14]{BASIS FOR $\Col B$}
\begin{itemize}
\item Let
\[
S=\{\vect b_1,\vect b_3,\vect b_5\}
=\left\{
\begin{bmatrix}1\\0\\0\\0\end{bmatrix},
\begin{bmatrix}0\\1\\0\\0\end{bmatrix},
\begin{bmatrix}0\\0\\1\\0\end{bmatrix}
\right\}.
\]
\item Since $\vect b_1\ne0$ and no vector in $S$ is a linear combination of the vectors that precede it, $S$ is linearly independent. (Theorem 4).
\item Thus $S$ is a basis for $\Col B$.
\end{itemize}
\end{frame}

% Source slide 47
\begin{frame}[shrink=14]{BASES FOR $\Nul A$ AND $\Col A$}
\begin{itemize}
\item \sourcehead{Theorem 6:} The pivot columns of a matrix $A$ form a basis for $\Col A$.
\item \textbf{Proof:} Let $B$ be the reduced echelon form of $A$.
\item The set of pivot columns of $B$ is linearly independent, for no vector in the set is a linear combination of the vectors that precede it.
\item Since $A$ is row equivalent to $B$, the pivot columns of $A$ are linearly independent as well, because any linear dependence relation among the columns of $A$ corresponds to a linear dependence relation among the columns of $B$.
\end{itemize}
\end{frame}

% Source slide 48
\begin{frame}[shrink=13]{BASES FOR $\Nul A$ AND $\Col A$}
\begin{itemize}
\item For this reason, every nonpivot column of $A$ is a linear combination of the pivot columns of $A$.
\item Thus the nonpivot columns of a may be discarded from the spanning set for $\Col A$, by the Spanning Set Theorem.
\item This leaves the pivot columns of $A$ as a basis for $\Col A$.
\item \textbf{Warning:} The pivot columns of a matrix $A$ are evident when $A$ has been reduced only to echelon form.
\item But, be careful to use the pivot columns of $A$ itself for the basis of $\Col A$.
\end{itemize}
\end{frame}

% Source slide 49
\begin{frame}[shrink=12]{BASES FOR $\Nul A$ AND $\Col A$}
\begin{itemize}
\item Row operations can change the column space of a matrix.
\item The columns of an echelon form $B$ of $A$ are often not in the column space of $A$.
\item \textbf{Two Views of a Basis}
\item When the Spanning Set Theorem is used, the deletion of vectors from a spanning set must stop when the set becomes linearly independent.
\item If an additional vector is deleted, it will not be a linear combination of the remaining vectors, and hence the smaller set will no longer span $V$.
\end{itemize}
\end{frame}

% Source slide 50
\begin{frame}{TWO VIEWS OF A BASIS}
\begin{itemize}
\item Thus a basis is a spanning set that is as small as possible.
\item A basis is also a linearly independent set that is as large as possible.
\item If $S$ is a basis for $V$, and if $S$ is enlarged by one vector---say, $\vect w$---from $V$, then the new set cannot be linearly independent, because $S$ spans $V$, and $\vect w$ is therefore a linear combination of the elements in $S$.
\end{itemize}
\end{frame}
